%% ================================================================
%%  SECTION 8 — Datasets
%% ================================================================
\section{Datasets}

Datasets have been the bottleneck. Western legal AI inherited decades of structured case law from repositories such as LexisNexis and PACER; Indian legal NLP has had to build its benchmarks from scratch, one annotation campaign at a time, typically depending on legal professionals whose time is expensive and scarce. We survey the seven principal datasets available as of 2025. They differ sharply in scale, court coverage, annotation depth, and public availability, and we believe these differences matter more than they are usually credited with: where the field claims to stand and where it actually stands diverge along precisely these axes.

\begin{table*}[t]
\centering
\caption{Primary Datasets for Indian Legal NLP Research (2021--2025)}
\label{tab:datasets}
\begin{tabular}{>{\raggedright\arraybackslash}p{2.4cm} >{\centering\arraybackslash}p{1.8cm} >{\raggedright\arraybackslash}p{1.8cm} >{\raggedright\arraybackslash}p{2.4cm} >{\raggedright\arraybackslash}p{4.2cm} >{\centering\arraybackslash}p{1.2cm}}
\toprule
\textbf{Dataset} & \textbf{Reference} & \textbf{Size} & \textbf{Court Coverage} & \textbf{Annotation Type} & \textbf{Avail.} \\
\midrule
ILDC & \cite{malik-etal-2021-ildc} & 35,000 cases & Supreme Court of India & Binary outcome labels, multi-sentence explanations & Public \\
\addlinespace
PredEx & \cite{nigam-etal-2024-predex} & 15,000+ cases & Supreme Court of India & Expert-annotated natural language explanations aligned to outcomes & Public \\
\addlinespace
NyayaAnumana & \cite{nigam-etal-2025-nyayaanumana} & 700,000+ cases & Supreme Court, High Courts, District Courts & Outcome labels, court-level metadata & Public \\
\addlinespace
NyayaRAG Corpus & \cite{nigam-etal-2025-nyayarag} & Subset of Indian judgments & Supreme Court and High Courts & Retrieval-annotated with relevant precedent links & Partial \\
\addlinespace
LegalEval (SemEval-2023 Task 6) & \cite{modi-etal-2023-semeval} & 9,380 sent. (NER); 14,864 sent. (roles) & Supreme Court and High Courts & NER labels (14 types), rhetorical role labels, statute tags & Public \\
\addlinespace
InLegalBERT Corpus & \cite{paul-etal-2023-inlegalbert} & 5.4M sentences & Supreme Court and High Courts & Pre-training corpus, no task-specific labels & Partial \\
\addlinespace
Indian Kanoon & -- & Millions of documents & All court levels & No annotation, raw text only & Public \\
\bottomrule
\end{tabular}
\end{table*}

The jump from ILDC's thirty-five thousand Supreme Court cases, sourced from Indian Kanoon and annotated with binary outcome labels alongside multi-sentence explanations~\cite{malik-etal-2021-ildc}, to NyayaAnumana's seven hundred thousand judgments, drawn from the Supreme Court, the High Courts, and the District Courts and shipped with outcome labels plus court-level metadata~\cite{nigam-etal-2025-nyayaanumana}, is not incremental. It is a qualitative shift that, for the first time, makes LLM-scale fine-tuning over Indian legal text genuinely feasible. Scale and depth, however, are not the same thing. NyayaAnumana provides outcome labels and court-level metadata at volume; it does not include the expert-written natural language explanations that \texttt{PredEx}~\cite{nigam-etal-2024-predex}, a fifteen-thousand-case Supreme Court corpus annotated by legal experts with explanations aligned to outcomes, supplies at a fraction of the scale. We would argue that any claim about ``state-of-the-art'' Indian legal judgment prediction needs to specify whether the model was trained with explanation supervision or without it, because the two settings produce fundamentally different systems and the literature has been sloppy about treating them as comparable.

District court data, meanwhile, is almost invisible across all annotated datasets, a glaring omission given that district courts handle the vast majority of India's pending caseload. The \texttt{InLegalBERT} pre-training corpus~\cite{paul-etal-2023-inlegalbert}, roughly 5.4 million sentences drawn from Supreme Court and High Court judgments and shipped without task-specific labels, together with the broader Indian Kanoon collection of millions of unannotated documents, supplies raw text at scale. With no task-specific annotations, however, their utility is limited to pre-training and unsupervised exploration.

Three gaps stand out, and none of the seven datasets we surveyed addresses any of them. First, there are no argument-level annotations linking claims to evidence, statutes, and precedents within a single judgment, which forces argument graph construction (Section~V) to operate in an unsupervised or weakly supervised mode. Second, there are no multilingual annotations that handle the Hindi--English or regional-language code-mixing pervasive in Indian judicial text, despite this code-mixing being routine rather than rare. Third, there are no counterfactual annotations of the kind a what-if analysis system (Section~VII) would need for either training or evaluation. Each of these gaps maps directly onto an open research direction we have flagged earlier in this paper. Closing them will require annotation campaigns that bring NLP researchers and legal domain experts together at a scale the community has not yet managed, and we are honestly not sure what the right organisational form for that collaboration looks like~\cite{modi-etal-2023-semeval}.

%% ================================================================
%%  SECTION 9 — Evaluation Metrics
%% ================================================================
\section{Evaluation Metrics}

How a system is measured matters at least as much as how it is built, and this is a point that legal AI papers tend to gloss over. Pick the wrong metric and you end up optimising for something the legal system does not actually care about, or worse, producing a model that posts impressive leaderboard numbers while failing on precisely the cases that matter most. The majority of legal AI studies report accuracy as their sole performance measure, without justifying that choice or acknowledging what it conceals~\cite{medvedeva-etal-2020-echr}. On imbalanced datasets, and Indian legal datasets are almost always imbalanced, accuracy is close to meaningless. We survey the metrics used across the retrieval, classification, and generation tasks covered in this paper, and where appropriate, we are not shy about flagging where metric choice has led the field astray.

\subsection{Retrieval Metrics}

\texttt{MAP} (Mean Average Precision) remains the default for ranked retrieval evaluation. It computes precision at each relevant document in the ranked list, averages across the list, and then averages again across queries, yielding a single number that captures both how many relevant documents were retrieved and how high they were ranked. \texttt{nDCG} (Normalised Discounted Cumulative Gain) complements MAP by applying a logarithmic discount to lower-ranked results, which aligns reasonably well with how lawyers actually search: they look at the first page, perhaps the second, and rarely venture further. Both BM25 evaluations~\cite{rosa-etal-2021-bm25} and dense retrieval experiments in COLIEE~\cite{althammer-etal-2021-dossier} report MAP and nDCG as primary measures.

The catch, and we believe this is underappreciated, is that both metrics require relevance judgments. In legal retrieval, relevance is not a surface-level property. Whether a 1987 Supreme Court ruling on tenant eviction is ``relevant'' to a 2022 high court dispute over landlord obligations is a question of legal interpretation, not of keyword overlap. Collecting those judgments at scale requires legal expertise, which means that evaluation itself is bottlenecked by exactly the same annotation scarcity that constrains training data, and we have not yet seen a methodology that genuinely works around this rather than around the edges of it.

\subsection{Classification, Generation, and Explainability Metrics}

Accuracy and macro-averaged F1 are the standard pair for binary and multi-class outcome prediction. On balanced data, accuracy is acceptable. Indian legal datasets, however, are rarely balanced. Dismissals dominate most corpora, which means a classifier that predicts ``dismissed'' every single time can post a respectable accuracy score without having learned anything at all~\cite{medvedeva-etal-2020-echr}. Macro-F1 corrects for this by weighting each class equally, and we would argue, plainly, that it should be the primary metric in every legal judgment prediction paper, not a secondary number buried in an appendix.

NER evaluation requires a finer granularity than aggregate F1 can provide. An aggregate score across all entity types can hide the fact that a model performs well on common categories such as dates and court names while failing almost completely on rare but legally critical ones such as witness names and provision numbers~\cite{kalamkar-etal-2022-ner}. Per-entity-type F1 is, in our view, the minimum acceptable standard; anything less is misleading.

Explanation evaluation is harder, because the property being measured, namely legal soundness, is not something an automated metric can assess. \texttt{ROUGE} and \texttt{BERTScore} capture surface overlap with reference text, which tells you whether the generated explanation reuses similar words. They tell you nothing about whether the underlying reasoning is correct. Expert Likert-scale assessments, in which legal professionals rate adequacy, accuracy, and completeness on structured scales, remain the only reliable approach. Both \texttt{INLegalLlama}~\cite{nigam-etal-2025-nyayaanumana} and \texttt{PredEx}~\cite{nigam-etal-2024-predex} adopted this methodology. It is expensive. It does not scale well. And it is, as far as we can tell, irreplaceable, which is a problem the community has not yet confronted honestly.
